\section{MCP Protocol}
\label{sec:mcp-protocol}

\subsection{Background}

The Model Context Protocol~\cite{anthropic2024mcp} is an open JSON-RPC~2.0
protocol, inspired by the Language Server Protocol~\cite{microsoft2016lsp},
designed for connecting LLM applications to external data sources and tools.
Like LSP, it runs over stdio with newline-delimited JSON messages. Unlike LSP
(which is editor-to-language-server), MCP is agent-to-tool-server: the
client is an LLM or AI agent, and the server exposes \emph{tools} --- named
callable operations with typed input schemas.

\subsection{Why MCP for code intelligence}

The tool-call pattern maps naturally to code intelligence queries. An AI agent
can call \texttt{symbol\_search} as naturally as it calls any other function.
The structured JSON input/output means the agent never needs to parse
grep-style output or interpret editor hover text. codetopo registers 27 tools
that cover the full range of structural queries an AI agent needs when
working on a codebase.

\subsection{Transport and framing}

\begin{itemize}
  \item Transport: stdio (one server process per client session).
  \item Framing: newline-delimited JSON (NDJSON); each message is a UTF-8
        JSON object terminated by \verb|\n|.
  \item Protocol version: \texttt{2024-11-05}\footnote{MCP protocol versions
        use calendar-date strings. \texttt{2024-11-05} is the version at which
        the core tool-call, prompts, and resources primitives were stabilized
        in the initial public specification.} (returned in the handshake
        \texttt{initialize} response).
\end{itemize}

\subsection{Connection lifecycle}

\begin{enumerate}
  \item Client sends \texttt{initialize} with protocol version and
        client capabilities.
  \item Server responds with \texttt{protocolVersion}, server info, and the
        full tool list.
  \item Client sends \texttt{initialized} notification.
  \item Client sends \texttt{tools/call} requests; server responds with results.
  \item Either side may send \texttt{shutdown} to end the session.
\end{enumerate}

\subsection{Staleness detection}

Each tool response includes a \texttt{\_meta} object. When the server detects
that the index may be out of date --- by checking the mtime of
\texttt{.git/HEAD} against the timestamp of the last index run --- it sets
\texttt{\_meta.stale = true}. Clients should treat stale responses as advisory:
the data is the last known good state, but a call to \texttt{reindex} is
recommended.

\subsection{QueryCache}

The \texttt{QueryCache} is an in-process prepared-statement cache keyed by
tool name. Preparing a SQLite statement compiles it to bytecode; caching avoids
re-compilation on every tool call. The cache is cleared (all statements
finalized) at the start of each re-index cycle so that schema changes take
effect immediately.

\subsection{Tool catalogue}

The 27 tools are organized into eight categories.

\subsubsection{Navigation}

\begin{description}
  \item[\texttt{dir\_list}] List files and directories under a path.
  \item[\texttt{file\_search}] Find files by name pattern across all roots.
  \item[\texttt{file\_summary}] Return symbol-annotated outline of a file.
  \item[\texttt{source\_at}] Read raw source content at a given path.
\end{description}

\subsubsection{Symbol lookup}

\begin{description}
  \item[\texttt{symbol\_search}] FTS5 prefix search over symbol names,
        qualified names, signatures, and docstrings. \texttt{kind='function'}
        matches both \texttt{function} and \texttt{method} nodes for
        polyglot parity.
  \item[\texttt{symbol\_list}] List all symbols of a given kind in a root.
  \item[\texttt{symbol\_get}] Fetch a single symbol by ID.
  \item[\texttt{symbol\_get\_batch}] Fetch multiple symbols by ID in one
        round-trip.
\end{description}

\subsubsection{Graph traversal}

\begin{description}
  \item[\texttt{callers\_approx}] Direct callers of a symbol (inbound
        \texttt{calls} edges). Named ``approx'' because call edges are
        AST-derived without full type resolution.
  \item[\texttt{callees\_approx}] Direct callees of a symbol (outbound
        \texttt{calls} edges), including unresolved name-only entries.
  \item[\texttt{references}] All reference sites for a symbol across all roots.
  \item[\texttt{context\_for}] Return the symbol context (enclosing function
        or class) for a given file location.
  \item[\texttt{entrypoints}] Identify symbols with no callers (potential
        entry points or public API).
\end{description}

\subsubsection{Impact analysis}

\begin{description}
  \item[\texttt{impact\_of}] Transitive BFS/DFS over the call graph from a
        seed node; returns all symbols reachable via \texttt{calls} edges.
  \item[\texttt{subgraph}] Extract the subgraph reachable from a set of seed
        nodes up to a given depth.
  \item[\texttt{shortest\_path}] Find the shortest call path between two
        symbols using BFS.
  \item[\texttt{dependency\_cluster}] Group symbols by mutual dependency into
        clusters.
  \item[\texttt{find\_implementations}] Find concrete implementations of an
        interface, trait, or abstract method.
\end{description}

Figure~\ref{fig:dependency-clusters} illustrates the partitioned output shape
of \texttt{dependency\_cluster}: symbols that share strong mutual dependencies
collapse into the same set.

\begin{figure}[htbp]
  \centering
  \tritonnext{width=0.68\textwidth}
\begin{triton}
unionfind 7
  title dependency clusters
  parent 0 0 1 3 3 5 5
\end{triton}
  \caption{A disjoint-set view of \texttt{dependency\_cluster} output, where
           each connected group is a coupled symbol cluster.}
  \label{fig:dependency-clusters}
\end{figure}

\subsubsection{Metadata}

\begin{description}
  \item[\texttt{repo\_stats}] Row counts, language distribution, and index
        size.
  \item[\texttt{server\_info}] Server version, protocol version, and
        supported languages.
  \item[\texttt{reindex}] Trigger an incremental or full re-index of all roots.
\end{description}

\subsubsection{File analysis}

\begin{description}
  \item[\texttt{file\_deps}] Return the import/include dependency graph for a
        file.
  \item[\texttt{method\_fields}] List the fields accessed by a method (field
        access edges).
\end{description}

\subsubsection{Workspace management}

\begin{description}
  \item[\texttt{workspace\_add}] Add a new root path to the workspace.
  \item[\texttt{workspace\_remove}] Remove a root and cascade-delete its data.
  \item[\texttt{workspace\_list}] List all registered roots with status and
        statistics.
\end{description}

\subsubsection{Full-text search}

\begin{description}
  \item[\texttt{code\_search}] Trigram full-text search over source lines
        using the \texttt{content\_fts} contentless FTS5 index~\cite{cox2012trigram}.
        Returns matching lines with file path and line number.
\end{description}
